\documentclass[10pt,a4paper]{amsart}
\usepackage[margin=19mm]{geometry}
\usepackage{amsmath,amssymb,mathtools}
\usepackage[scr=rsfs,cal=euler]{mathalfa}
\usepackage{xfrac}
\usepackage[unicode,hidelinks,bookmarks=false]{hyperref}
\newcommand{\ie}{i.e.\ }
\newcommand{\cf}{cf.\ }
\newcommand{\resp}{respectively\ }
\newcommand{\Subsection}{\subsection}
\providecommand{\nobreakdash}{\nobreak\hbox{-}}
\setlength{\emergencystretch}{2em}
\multlinegap0pt
\usepackage{amssymb}
\usepackage[normalem]{ulem}
\usepackage{mathtools}
\usepackage{xcolor}
\newtheorem{prop}{Proposition}
\newcommand{\R}{\mathbb{R}}
\newcommand{\T}{{\mathbb{T}}}
\title{full title}
\author{Alessio Basti, Fabio Camilli}
\date{Source: arXiv:2512.24051v2 (CC BY 4.0)}
\begin{document}
\maketitle

\noindent\emph{Source and licence.} full title. Version arXiv:2512.24051v2 (CC BY 4.0), distributed by arXiv under the Creative Commons Attribution 4.0 International licence, http://creativecommons.org/licenses/by/4.0/, as recorded in the arXiv licence metadata for this exact version and shown on the abstract page https://arxiv.org/abs/2512.24051v2. Full licence notice: \texttt{licenses/arxiv-2512.24051-CC-BY-4.0.txt}.\par\smallskip

\noindent\textit{Editorial scope.} Counted unit: the article's prop prop:properties\_scheme (source0.tex lines 339--364). The article's complete original proof of it is delivered with it (source0.tex lines 938--1158, one \texttt{proof} environment of 221 lines, which the article places away from the statement). No mathematics is written, repaired or re-derived: the statement and the proof are the article's own lines, in the article's own notation, and no retained line is substituted or re-flowed.  Retained uncounted as the source-local context the proof needs: lines 293--298.  Nothing was omitted from any retained span, so the package carries no omission record.  No retained line contains a citation, so the wrapper carries no bibliography and no bibliography entry.  No retained line contains a figure environment, an image-inclusion command or drawing code, so no image is delivered and the package declares no asset dependency.  Editorial formatting only, in addition: an AMS \texttt{amsart} preamble that restates the article's own declarations and macros used by the retained lines, namely the environments \texttt{prop} on separate counters, together with the standard packages those lines need.  The retained environments print in this document as Proposition 1 in order of appearance, whereas the article prints the same items with the numbering of its own full document.  The complete native source of the article as distributed in the arXiv e-print for this exact version is bundled unchanged as \texttt{sources/191847/source0.tex}.
\begin{equation}\label{HJ_app_d}
\begin{cases}
u^\varepsilon_{n+1}(x)-u^\varepsilon_n(x) + \Delta t\, F\big(-\delta_{h} u^\varepsilon_n(x),\ \delta_{-h} u^\varepsilon_n(x)\big)=0, & x\in\mathbb{T}^d,\ n=0,\dots,N-1,\\[4pt]
u^\varepsilon_0(x)=g(x), & x\in\mathbb{T}^d,
\end{cases}
\end{equation}

\begin{prop}[Stability and semiconcavity]\label{prop:properties_scheme}
Let \(\{u^\varepsilon_n\}_{n=0}^N\) be the solution of \eqref{HJ_app_d}. 
Under assumptions \textup{(H1)--(H2)}, \textup{(G1)} (with semiconcavity constant \(C_{conc}\)), and \textup{(F1)--(F3)}, the following properties hold for every \(n=0,\dots,N\):
\begin{enumerate}
  \item[\textup{(i)}] Lipschitz bounds (discrete and continuous).
  \[
    \|\delta_h u^\varepsilon_n\|_{L^\infty(\T^d)} \le \|\delta_h g\|_{L^\infty(\T^d)},
    \qquad
    \|Du^\varepsilon_n\|_{L^\infty(\T^d)} \le \|Dg\|_{L^\infty(\T^d)},
  \]
  where the gradient \(Du^\varepsilon_n\) is a.e.
  \item[\textup{(ii)}] Stability.
  \[
    \|u^\varepsilon_n\|_{L^\infty(\T^d)} \le \|g\|_{L^\infty(\T^d)} + n\Delta t\,|H(0)|.
  \]
  \item[\textup{(iii)}] Semiconcavity. For every \(x\in\T^d\) and every \(k\in\T^d\setminus\{0\}\) define
  \[
    \Delta_k u^\varepsilon_n(x):=u^\varepsilon_n(x+k)+u^\varepsilon_n(x-k)-2u^\varepsilon_n(x).
  \]
  Then
  \begin{equation}\label{semiconc_solutions_finite_diff}
    \sup_{x\in\T^d,\; k\in \T^d\setminus\{0\}}\frac{\Delta_k u^\varepsilon_n(x)}{|k|^2} \le C_{conc},
  \end{equation}
  i.e. the semiconcavity constant of \(u^\varepsilon_n\) is bounded by that of the initial datum \(g\).
\end{enumerate}
\end{prop}

\begin{proof}[Proof of Prop. \ref{prop:properties_scheme}]	

\noindent\textbf{(i)}
Let
\[
w_n^{(i)}(x):=\delta_h^{(i)}u_n^\varepsilon(x),\qquad
R_h:=\|\delta_h g\|_{L^\infty(\T^d)}
=\max_{1\le i\le d}\|\delta_h^{(i)}g\|_{L^\infty(\T^d)} .
\]
We first prove the discrete Lipschitz estimate by induction on \(n\). The base case
\(n=0\) is immediate. Assume the claim holds at time \(n\). Fix a direction \(i\).
By the induction hypothesis,
\[
u_n^\varepsilon(x+h e_i)-u_n^\varepsilon(x)
=h\,w_n^{(i)}(x)\le hR_h .
\]
Therefore
\[
u_n^\varepsilon(x+h e_i)\le u_n^\varepsilon(x)+hR_h .
\]
Consider the discrete evolution operator
\[
(\mathcal G_h v)(y):=
v(y)-\Delta t\,F\big(-\delta_h v(y),\delta_{-h}v(y)\big),
\]
so that \(u_{n+1}^\varepsilon=\mathcal G_h u_n^\varepsilon\). By the monotonicity of
\(\mathcal G_h\), and since adding a constant leaves the finite differences unchanged, we get
\[
\begin{aligned}
u_{n+1}^\varepsilon(x+h e_i)
&=(\mathcal G_h u_n^\varepsilon)(x+h e_i)  \\
&=\mathcal G_h\big(u_n^\varepsilon(\cdot+h e_i)\big)(x) \\
&\le \mathcal G_h\big(u_n^\varepsilon+hR_h\big)(x) \\
&=(\mathcal G_h u_n^\varepsilon)(x)+hR_h \\
&=u_{n+1}^\varepsilon(x)+hR_h .
\end{aligned}
\]
Equivalently,
\[
\delta_h^{(i)}u_{n+1}^\varepsilon(x)\le R_h .
\]
Applying the same argument to the inequality
\[
u_n^\varepsilon(x)\le u_n^\varepsilon(x+h e_i)+hR_h
\]
gives
\[
\delta_h^{(i)}u_{n+1}^\varepsilon(x)\ge -R_h .
\]
Hence
\[
\max_{1\le i\le d}
\|\delta_h^{(i)}u_{n+1}^\varepsilon\|_{L^\infty(\T^d)}
\le R_h .
\]
By induction,
\[
\|\delta_h u_n^\varepsilon\|_{L^\infty(\T^d)}
\le \|\delta_h g\|_{L^\infty(\T^d)}
\qquad\text{for every }n=0,\ldots,N .
\]
The same estimate for \(\delta_{-h}u_n^\varepsilon\) follows by translation.
It remains to prove the continuous Lipschitz estimate. Set
\[
L_g:=\|Dg\|_{L^\infty(\T^d)} .
\]
For \(z\in\R^d\), let \(T_zv(x):=v(x+z)\). Since \(g\) is Lipschitz,
\[
T_zg(x)=g(x+z)\le g(x)+L_g|z|
\qquad\text{for all }x\in\T^d,\ z\in\R^d .
\]
We prove by induction that
\[
u_n^\varepsilon(x+z)\le u_n^\varepsilon(x)+L_g|z|
\qquad\text{for all }x\in\T^d,\ z\in\R^d .
\]
The case \(n=0\) is the previous inequality for \(g\). Assume it holds at time \(n\).
Using the translation invariance, monotonicity, and constant-preserving property of
\(\mathcal G_h\), we obtain
\[
\begin{aligned}
u_{n+1}^\varepsilon(x+z)
&=(\mathcal G_h u_n^\varepsilon)(x+z) \\
&=\mathcal G_h(T_z u_n^\varepsilon)(x) \\
&\le \mathcal G_h(u_n^\varepsilon+L_g|z|)(x) \\
&=(\mathcal G_h u_n^\varepsilon)(x)+L_g|z| \\
&=u_{n+1}^\varepsilon(x)+L_g|z| .
\end{aligned}
\]
Replacing \(z\) by \(-z\) gives
\[
u_{n+1}^\varepsilon(x)-u_{n+1}^\varepsilon(x+z)\le L_g|z| .
\]
Therefore
\[
|u_{n+1}^\varepsilon(x+z)-u_{n+1}^\varepsilon(x)|
\le L_g|z|
\qquad\text{for all }x,z .
\]
Thus \(u_{n+1}^\varepsilon\) is \(L_g\)-Lipschitz. By induction,
\[
\operatorname{Lip}(u_n^\varepsilon)\le \operatorname{Lip}(g)
=\|Dg\|_{L^\infty(\T^d)} .
\]
Since \(u_n^\varepsilon\) is Lipschitz, it is differentiable a.e. by Rademacher's theorem, and
\[
\|D u_n^\varepsilon\|_{L^\infty(\T^d)}
\le \|Dg\|_{L^\infty(\T^d)} .
\]

		\noindent\textbf{(ii)}
Let
\(
w_n(x):=u_n^\varepsilon(x)+n\Delta t\,H(0),
\)
so that \(\delta_{\pm h}w_n=\delta_{\pm h}u_n^\varepsilon\). By the discrete Lipschitz bound proved in part (1) there exists \(R>0\) such that for every \(n\) and every coordinate \(i\)
\(
\sup_{x\in\T^d}|\delta^{(i)}_{\pm h} w_n(x)| \le R.
\)

From the scheme \eqref{HJ_app_d} we have pointwise for every \(x\in\T^d\)
\[
w_{n+1}(x)-w_n(x)
= -\Delta t\Big(F\big(-\delta_h w_n(x),\delta_{-h} w_n(x)\big)-F(0,0)\Big),
\]

Apply the multivariate mean value theorem to \(F:\R^{2d}\to\R\) on the segment joining \((0,0)\) and \(\big(-\delta_h w_n(x),\delta_{-h} w_n(x)\big)\). There exists \(\xi(x)\in\R^{2d}\) on that segment such that
\[
\begin{aligned}
F\big(-\delta_h w_n(x),\delta_{-h} w_n(x)\big)-F(0,0)
&= DF(\xi(x))\cdot\big(-\delta_h w_n(x),\delta_{-h} w_n(x)\big)\\
\end{aligned}
\]

Substituting into the equation of interest, yields the identity
\[
w_{n+1}(x)-w_n(x)
= -\frac{\Delta t}{h}\sum_{i=1}^d\Big[ F_{p_i}(\xi(x))\big(-w_n(x+h e_i)+w_n(x)\big)+ F_{q_i}(\xi(x))\big(w_n(x)-w_n(x-h e_i)\big)\Big],
\]
Therefore,
\[
\begin{aligned}
w_{n+1}(x)
&= \frac{\Delta t}{h}\sum_{i=1}^d F_{p_i}(\xi(x))\,w_n(x+h e_i)
+ \Big(1-\frac{\Delta t}{h}\sum_{i=1}^d\big(F_{p_i}(\xi(x))+F_{q_i}(\xi(x))\big)\Big) w_n(x)\\
&\qquad + \frac{\Delta t}{h}\sum_{i=1}^d F_{q_i}(\xi(x))\,w_n(x-h e_i).
\end{aligned}
\]
Since \(F_{p_i},F_{q_i}\ge0\), and the CFL condition holds, the three coefficients are nonnegative and sum to \(1\). It follows that, for each fixed \(x\),
	$$
		w_{n+1}(x)\le \max\{w_n(x),w_n(x\pm he_1),...,w_n(x\pm he_d)\}\le \sup_y w_n(y),
	$$
	hence \(\sup_x w_{n+1}(x)\le \sup_x w_n(x)\). 
	Similarly \(w_{n+1}(x)\ge \min\{w_n(x),w_n(x\pm he_1),...,w_n(x\pm he_d)\}\ge \inf_y w_n(y)\). Iterating we get \(\|w_n\|_{L^\infty(\T^d)}\le \|w_0\|_{L^\infty(\T^d)}=\|g\|_{L^\infty(\T^d)}\), and therefore
\[
\|u_n^\varepsilon\|_{L^\infty(\T^d)}
\le \|w_n\|_{L^\infty(\T^d)} + n\Delta t\,|H(0)|
\le \|g\|_{L^\infty(\T^d)} + n\Delta t\,|H(0)|.
\]


	\medskip
	\noindent\textbf{(iii)}
	Let \(w_n(x):=\Delta_k u^\varepsilon_n(x)/|k|^2\), that is, the quantity we want to estimate \ref{semiconc_solutions_finite_diff}.
	Introducing
	\(
	A_n(y):=\big(-\delta_h u^\varepsilon_n(y),\; \delta_{-h}u^\varepsilon_n(y)\big)
	\), evaluating the scheme \eqref{HJ_app_d} at the points \(y=x-k,x,x+k\), and suitably combining the three relations, we get
	\begin{equation}\label{eq:sn_evol}
		w_{n+1}(x)-w_n(x)
		= -\frac{\Delta t}{|k|^2}\Big( F\big(A_n(x+k)\big)-2F\big(A_n(x)\big)+F\big(A_n(x-k)\big)\Big).
	\end{equation}
	
	Since \(F\in C^2(\mathbb R^{2d})\) is convex, we have for any vectors \(X,Y,Z\in\mathbb R^{2d}\) \(F(X)-F(Y)\ge D F(Y)\cdot (X-Y)\) and \(F(Z)-F(Y)\ge D F(Y)\cdot (Z-Y)\). Therefore,
	\[
	F(X)-2F(Y)+F(Z)\ge D F(Y)\cdot (X-2Y+Z).
	\]
	Applying this inequality with
	\(X=A_n(x+k),\,Y=A_n(x),\,Z=A_n(x-k)\) in \eqref{eq:sn_evol} gives the estimate
	\[
		w_{n+1}(x)-w_n(x)
		\le -\frac{\Delta t}{|k|^2}\,D F\big(A_n(x)\big)\cdot\big(A_n(x+k)-2A_n(x)+A_n(x-k)\big)\]
Using \(\Delta_k u_n^\varepsilon=|k|^2 w_n\), and \(
\Delta_k\big(\delta^{(i)}_h u_n^\varepsilon\big)(x)
= \delta^{(i)}_h\big(\Delta_k u_n^\varepsilon\big)(x)
\), we get
\[
\Delta_k\big(\delta^{(i)}_h u_n^\varepsilon\big)(x)
=|k|^2\delta^{(i)}_h w_n(x),\qquad
\Delta_k\big(\delta^{(i)}_{-h} u_n^\varepsilon\big)(x)=|k|^2\delta^{(i)}_{-h} w_n(x),
\]
so that
\[
A_n(x+k)-2A_n(x)+A_n(x-k)=\Delta_k A_n(x)=|k|^2\big(-\delta_h w_n,\;\delta_{-h} w_n\big)(x).
\]
Therefore,
\[
w_{n+1}(x)-w_n(x)
\le -\Delta t\sum_{i=1}^d\big(-F_{p_i}(A_n(x))\,\delta^{(i)}_h w_n(x)
+F_{q_i}(A_n(x))\,\delta^{(i)}_{-h} w_n(x)\big).
\]
Writing \(\delta^{(i)}_{\pm h}w_n\) explicitly gives
\[
w_{n+1}(x)
\le w_n(x) -\frac{\Delta t}{h}\sum_{i=1}^d\Big(-F_{p_i}w_n(x+h e_i)+(F_{p_i}+F_{q_i})w_n(x)-F_{q_i}w_n(x-h e_i)\Big)
\]
\[
=w_n(x)\bigg(1 -\frac{\Delta t}{h}\sum_{i=1}^d(F_{p_i}+F_{q_i})\bigg)+\frac{\Delta t}{h}\sum_{i=1}^dF_{p_i}w_n(x+h e_i) +\frac{\Delta t}{h}\sum_{i=1}^dF_{q_i}w_n(x-h e_i).
\]
Since the scheme satisfies the monotonicity and CFL conditions, then the coefficients in the previous inequality are nonnegative and sum to $1$. Hence,
\[
w_{n+1}(x)\le\max\{w_n(x),\,w_n(x\pm h e_i)\}\le\sup_{y} w_n(y).
\]
Taking the supremum in $x$ yields $\sup_x w_{n+1}(x)\le\sup_x w_n(x)$, therefore the sequence $\{\sup_x w_n\}_n$ is nonincreasing. Consequently
\[
\sup_{x\in\T^d,\;k\in\T^d\setminus\{0\}}\frac{\Delta_k u_n^\varepsilon(x)}{|k|^2}
\le \sup_{x\in\T^d,\;k\in\T^d\setminus\{0\}} w_0(x)
= \sup_{x\in\T^d,\;k\in\T^d\setminus\{0\}}\frac{\Delta_k g(x)}{|k|^2}
\le C_{\mathrm{conc}}.
\]
\end{proof}

\end{document}
